% TOP_PROBLEM: {"id": 10100001, "title": "No percolation at the critical point on $\\mathbb{Z}^d$", "category_id": 19, "category": {"id": 19, "name": "probability", "display_name": "Probability", "description": "Problems involving probability theory, stochastic processes, and random structures.", "slug": "probability", "order_index": 19, "created_at": "2026-07-31T15:26:25.671Z"}, "status": "partially_solved"}
% FIRST_POSTED: null
% ATTEMPT_STATUS: unresolved
% Research continuation compiled from the conversation, 27 September 2026.
% Root statement and all substantive progress preserved; final audit leaves the problem unresolved.
\documentclass[11pt]{article}
\usepackage[margin=25mm]{geometry}
\usepackage{fontspec}
\setmainfont{Latin Modern Roman}
\usepackage{amsmath,amssymb,mathtools}
\usepackage{unicode-math}
\setmathfont{Latin Modern Math}
\usepackage{xurl,hyperref}
\hypersetup{colorlinks=true,urlcolor=blue}
\setcounter{secnumdepth}{0}
\setlength{\parindent}{0pt}
\setlength{\parskip}{5pt}
\setlength{\emergencystretch}{3em}
\newcommand{\one}{\mathbf 1}
\begin{document}

\section*{\texorpdfstring{No percolation at the critical point on $\mathbb{Z}^d$}{No percolation at the critical point on \$\textbackslash{}mathbb\{Z\}\textasciicircum{}d\$}}\label{critical-percolation}\label{10100001}

\subsection{Definitions and mathematical statement}
Consider nearest-neighbour Bernoulli bond percolation on \(\mathbb Z^d\). Each edge is independently open with probability \(p\in[0,1]\). Let
\[
\theta(p)=\mathbb P_p(0\leftrightarrow\infty),
\qquad
p_c=\inf\{p:\theta(p)>0\}.
\]
The conjecture is
\[
\boxed{\theta(p_c)=0.}
\]
Equivalently, at criticality there is almost surely no infinite open cluster.

This is often called the \(\theta(p_c)=0\) conjecture, the critical percolation conjecture, the conjecture that there is no infinite cluster at criticality, or the continuity conjecture for the percolation probability at \(p_c\). There is no single universally fixed eponymous name.

The discussion below specializes to nearest-neighbour Bernoulli bond percolation on \(\mathbb Z^d\). A status checkpoint cited during the conversation was that the conjecture is known in \(d=2\), known in sufficiently high dimension (in particular \(d\ge 11\) in the cited 2026 source), and unresolved in dimensions \(3\le d\le10\); see \href{https://arxiv.org/html/2608.23661v1}{[S1]}.

\subsection{Short English statement}
At the exact threshold where percolation first becomes possible, is the probability that the origin belongs to an infinite open cluster still zero?

\subsection{Research attempt}

\paragraph{Scope and outcome checkpoint.}
The aim was to obtain a full proof, not merely a literature summary. Several routes were developed substantially: proportional local uniqueness, robustness under thinning, a two-copy switching argument, a weighted contact lemma, a line-gate criterion, a finite-cluster radius-tail criterion, a block-Bessel inequality, and a decision-tree deletion inequality. The final audit uncovered a decisive conditioning gap in the proposed shell application. Consequently, the conjecture remains unresolved in this writeup. The valid reductions and auxiliary inequalities are retained below, together with the precise point where the attempted proof fails.

\paragraph{1. First reduction: proportional local uniqueness.}
For \(r>0\), write
\[
B_r=[-r,r]^d\cap\mathbb Z^d.
\]
For \(K>1\), let \(D_{L,K}\) be the event that, in the graph induced by \(B_{KL}\), there exist two distinct open clusters both meeting \(B_L\) and \(\partial B_{KL}\). The two clusters inside \(B_{KL}\) may still be pieces of the same global infinite cluster reconnecting outside \(B_{KL}\).

The proposed proportional local uniqueness property was
\[
\tag{PLU}\label{eq:plu}
\lim_{K\to\infty}\limsup_{L\to\infty}\mathbb P_p(D_{L,K})=0.
\]

\textbf{Conditional proposition.}
If \(\theta(p)>0\) and \(\mathrm{PLU}(p)\) hold, then \(p>p_c\).

\emph{Argument.}
If \(\theta(p)>0\), sufficiently large boxes are very likely to meet the infinite cluster. Under \eqref{eq:plu}, for suitable fixed \(K\) and large \(L\), a finite collection of translated \(KL\)-boxes can be declared good so that the central region connects to prescribed port regions and macroscopic crossing clusters within the enlarged box consolidate into a distinguished crossing structure. Neighbouring good blocks can then be arranged to have crossing clusters that join.

The good-block indicators form a finite-range dependent process. For marginal probability sufficiently close to one, the Liggett--Schonmann--Stacey theorem gives domination of a supercritical independent percolation process; see \href{https://projecteuclid.org/journals/annals-of-probability/volume-25/issue-1/Domination-by-product-measures/10.1214/aop/1024404279.pdf}{[S3]}.

For fixed \(K,L\), goodness is a finite event, so its probability is a polynomial in the edge density \(q\), hence continuous in \(q\). Once it lies strictly above the coarse percolation threshold at \(p\), it remains above it for some \(q<p\). Thus \(q\ge p_c\), proving \(p>p_c\).

Hence
\[
\boxed{\theta(p)>0+\mathrm{PLU}(p)\Longrightarrow p>p_c.}
\]
Applying this at \(p=p_c\) would prove \(\theta(p_c)=0\).

The known qualitative statement has the quantifier pattern
\[
\forall m\,\forall\varepsilon\,\exists M=M(m,\varepsilon),
\]
whereas the needed proportional version asks for
\[
\forall\varepsilon\,\exists K<\infty\,\forall L\gg1,\qquad M=KL.
\]
The dependence \(M/m\) is therefore the first precise obstruction.

\paragraph{2. Thinning reformulation.}
Assume for contradiction that \(\theta(p_c)>0\), and independently retain each \(p_c\)-open edge with probability \(r<1\). The retained configuration is Bernoulli percolation with parameter \(rp_c<p_c\), so it cannot percolate. Hence a hypothetical critical infinite cluster \(\mathcal C_\infty\) would satisfy
\[
\boxed{p_c^{\mathrm{bond}}(\mathcal C_\infty)=1}
\]
under independent thinning: deleting any positive density of its edges destroys every infinite component.

Thus a critical infinite cluster would be positive-density but maximally fragile. Density and uniqueness alone do not exclude this for arbitrary subgraphs, so the missing ingredient must use the special Bernoulli geometry.

\paragraph{3. First two-copy switching attempt.}
Kozma and Nitzan proposed a finite-graph inequality
\[
\tag{KN}\label{eq:kn}
\mathbb P(0\leftrightarrow b)
\ge
\mathbb P(0\leftrightarrow A)
\min_{a\in A}\mathbb P(a\leftrightarrow b),
\]
or a suitable high-probability variant, as a route to \(\theta(p_c)=0\); see \href{https://arxiv.org/html/2401.12397v1}{[S2]}.

Take two independent configurations \(\omega,\eta\). On \(0\leftrightarrow A\) in \(\omega\), select
\[
a(\omega)\in A\cap\mathcal C_\omega(0).
\]
Then
\[
\mathbb P\bigl(
0\leftrightarrow A\text{ in }\omega,\,
a(\omega)\leftrightarrow b\text{ in }\eta
\bigr)
\ge
\mathbb P(0\leftrightarrow A)
\min_{a\in A}\mathbb P(a\leftrightarrow b).
\]
The natural plan is to splice the \(\omega\)-path to \(a(\omega)\) with an \(\eta\)-path onward to \(b\). For a fixed edge set, swapping the two independent copies preserves law, but an adaptively chosen switching set need not. A one-edge warning example is that adaptively placing an open bit in the first coordinate whenever either of two Bernoulli-\(p\) copies is open changes its open probability to
\[
1-(1-p)^2=2p-p^2>p.
\]
Thus the naive adaptive switching proof is invalid.

\paragraph{4. Finite-cluster-tail route.}
Cerf's August 2026 preprint gives, for Bernoulli site percolation on \(\mathbb Z^d\), a critical alternative under which a hypothetical infinite cluster at \(p_c\) forces failure of stretched-exponential finite-cluster tails; see \href{https://arxiv.org/html/2608.23661v1}{[S1]}.

This suggests trying to prove from \(\theta(p)>0\) alone that, for some \(c,\alpha>0\),
\[
\mathbb P_p\bigl(n\le |C(0)|<\infty\bigr)\le e^{-c n^\alpha}.
\]
Such an estimate at \(p=p_c\) would contradict the alternative. Classical estimates of this type are established in the strictly supercritical regime \(p>p_c\), however; deriving them merely from \(\theta(p)>0\) is another form of the same robustness problem.

\paragraph{5. Weighted-contact gluing lemma.}
Let \(G\) be a finite graph, \(s,t\in G\), and \(A\subseteq G\). Give \(A\) nonnegative weights \(w_a\), and put
\[
X=\sum_{a\in A}w_a\,\one_{\{s\leftrightarrow a\}}.
\]

\textbf{Lemma 5.1 (contact-mass gluing).}
For every \(r>0\),
\[
\boxed{
\mathbb P(s\leftrightarrow t)
\ge
\mathbb P(X\ge r)
-\frac1r\sum_{a\in A}
w_a\,\mathbb P(s\leftrightarrow a)\mathbb P(a\nleftrightarrow t).
}
\tag{5.1}\label{eq:contactmass}
\]
If \(\mathbb P(a\leftrightarrow t)\ge1-\delta\) for every \(a\in A\), then
\[
\mathbb P(s\leftrightarrow t)
\ge
\mathbb P(X\ge r)-\frac{\delta\,\mathbb EX}{r}.
\]
If moreover \(\sum_aw_a=1\), then
\[
\boxed{
\mathbb P(s\leftrightarrow t)
\ge
\mathbb P(X\ge r)-\frac{\delta}{r}.
}
\tag{5.2}
\]

\emph{Proof.}
Let \(F=\{s\nleftrightarrow t\}\). On \(F\), \(s\leftrightarrow a\) implies \(a\nleftrightarrow t\), so
\[
X\one_F
\le
\sum_{a\in A}w_a\one_{\{s\leftrightarrow a,\ a\nleftrightarrow t\}}.
\]
By increasing--decreasing Harris--FKG,
\[
\mathbb P(s\leftrightarrow a,\ a\nleftrightarrow t)
\le
\mathbb P(s\leftrightarrow a)\mathbb P(a\nleftrightarrow t).
\]
Therefore
\[
\mathbb E[X\one_F]
\le
\sum_{a\in A}w_a\mathbb P(s\leftrightarrow a)\mathbb P(a\nleftrightarrow t).
\]
Since \(X\ge r\) on \(F\cap\{X\ge r\}\),
\[
r\,\mathbb P(F,X\ge r)\le\mathbb E[X\one_F].
\]
Subtracting the bad part of \(\{X\ge r\}\) proves the result. \hfill\(\square\)

Thus it would suffice to show that the source reaches a non-negligible weighted mass of target-good shell components. Existing shell arguments guarantee contact with some good component but do not control its weight.

\paragraph{6. Parallel-line gate inequality.}
Let
\[
x_i=i e_1,\qquad y_i=i e_1+e_2,\qquad e_i=\{x_i,y_i\}
\]
for \(1\le i\le n\). Define
\[
A_n=\{x_i\nleftrightarrow\infty\text{ for every }1\le i\le n\},
\qquad
M_n=\sum_{i=1}^n\one_{\{y_i\leftrightarrow\infty\}}.
\]

\textbf{Lemma 6.1 (parallel-line gates).}
For every integer \(k\ge1\),
\[
\boxed{
\mathbb P_p(A_n)
\le
\mathbb P_p(M_n<k)+(1-p)^k.
}
\tag{6.1}\label{eq:gates}
\]

\emph{Proof.}
Condition on all edges except \(E_n=\{e_1,\ldots,e_n\}\). With the gate edges deleted, let
\[
R=\{i:y_i\leftrightarrow\infty\}.
\]
If \(i\in R\) and \(e_i\) is open, then \(x_i\leftrightarrow\infty\), so \(A_n\) forces all \(e_i\) with \(i\in R\) to be closed. On \(A_n\), any \(y_i\) infinite in the completed configuration must already lie in \(R\), otherwise its infinite path would traverse some gate \(e_j\) and make \(x_j\) infinite. Hence on \(A_n\), \(M_n=|R|\). Conditional on all nongate edges,
\[
\mathbb P_p(A_n,M_n\ge k\mid\omega_{E_n^c})\le(1-p)^k.
\]
Averaging proves \eqref{eq:gates}. \hfill\(\square\)

Choose
\[
k_n=\lceil C\log n\rceil,
\qquad
C>\frac{2}{|\log(1-p_c)|}.
\]
Then \(\sum_n(1-p_c)^{k_n}<\infty\). Therefore it would suffice to prove
\[
\boxed{
\sum_{n=1}^\infty
\mathbb P_{p_c}(M_n<C\log n)<\infty.
}
\tag{6.2}\label{eq:linecriterion}
\]
Indeed, \(\sum_n\mathbb P_{p_c}(A_n)<\infty\), implying finite mean distance to the next infinite-cluster point on a parallel line, which the Kozma--Nitzan discussion ties to the Barsky--Grimmett--Newman half-space theorem.

Ergodicity only gives \(M_n/n\to\theta(p)\) almost surely; it gives no summable convergence rate.

\paragraph{7. Finite-cluster radius-tail criterion.}
Define
\[
q_p(R)=
\mathbb P_p\bigl(
0\leftrightarrow\partial B_R,\,
0\nleftrightarrow\infty
\bigr).
\tag{7.1}\label{eq:truncatedarm}
\]

\textbf{Proposition 7.1.}
Suppose \(\theta(p)=\vartheta>0\). There are \(c_1,c_2>0\), depending only on \(\vartheta\), such that whenever \(R\le n/10\) and
\[
k\le c_1\vartheta\,\frac nR,
\]
one has
\[
\boxed{
\mathbb P_p(M_n<k)
\le
\exp\left(-c_2\vartheta\,\frac nR\right)
+\frac4{\vartheta}q_p(R).
}
\tag{7.2}\label{eq:radiustail}
\]

\emph{Proof.}
Choose \(m\asymp n/R\) points \(z_1,\ldots,z_m\) among \(y_1,\ldots,y_n\), separated by more than \(2R+2\), so that the edge sets of \(B_R(z_j)\) are disjoint. Let
\[
Z_j=\one_{\{z_j\leftrightarrow\partial B_R(z_j)\text{ inside }B_R(z_j)\}}.
\]
The \(Z_j\) are independent and
\[
\mathbb P(Z_j=1)=\vartheta+q_p(R)\ge\vartheta.
\]
A Chernoff bound gives
\[
\mathbb P\left(\sum_{j=1}^mZ_j<\frac{\vartheta m}{2}\right)
\le e^{-c\vartheta m}.
\]
Let \(F_j\) indicate that \(Z_j=1\) but \(z_j\nleftrightarrow\infty\). Since \(\mathbb EF_j=q_p(R)\), Markov gives
\[
\mathbb P\left(\sum_{j=1}^mF_j\ge\frac{\vartheta m}{4}\right)
\le\frac{4q_p(R)}{\vartheta}.
\]
Outside these two bad events, at least \(\vartheta m/4\) sampled points belong to an infinite cluster. Since \(m\asymp n/R\), the estimate follows after adjusting constants. \hfill\(\square\)

Taking
\[
k_n=\lceil C\log n\rceil,
\qquad
R_n=\left\lfloor\frac{n}{K\log n}\right\rfloor
\]
gives the sufficient criterion
\[
\boxed{
\sum_{n=2}^\infty
q_{p_c}\left(\left\lfloor\frac{n}{K\log n}\right\rfloor\right)<\infty
\Longrightarrow
\theta(p_c)=0.
}
\tag{7.3}\label{eq:summableq}
\]
In particular, a bound
\[
q_{p_c}(R)=O(R^{-1-\alpha}),\qquad \alpha>0,
\tag{7.4}
\]
or a comparable logarithmically improved \(1/R\) bound such as
\[
q_{p_c}(R)
=
O\left(\frac1{R(\log R)^{2+\alpha}}\right),
\tag{7.5}
\]
would suffice.

The attempted BK/Reimer amplification of \(q_p(R)\to0\) did not close: several sampled arms may lie in the same finite cluster; finiteness is not an increasing event; separation certificates can overlap; and buffering merely moves the unresolved tail to a larger scale.

\paragraph{8. Block-Bessel aggregation of rare connector seeds.}
Let \(B_1,\ldots,B_N\) be pairwise edge-disjoint blocks in a finite graph. Put
\[
F_i=\{\text{all edges of }B_i\text{ are open}\},
\quad
q_i=\mathbb P(F_i),
\quad
w_i=\frac{q_i}{1-q_i},
\quad
Q=\sum_iw_i.
\]
Suppose forcing \(B_i\) open deterministically connects \(s\) to \(a_i\), and let
\[
\tau_i=\mathbb P(a_i\leftrightarrow t),
\qquad
\overline\tau=\frac{\sum_iw_i\tau_i}{Q}.
\]

\textbf{Lemma 8.1 (connector-block Bessel inequality).}
\[
\boxed{
\mathbb P(s\leftrightarrow t)
\ge
\overline\tau-\frac1{2\sqrt Q}.
}
\tag{8.1}\label{eq:blockbessel}
\]

\emph{Proof.}
Let \(C=\{s\leftrightarrow t\}\), \(c=\mathbb P(C)\), and
\[
D_i=\frac{\one_{F_i}}{q_i}-1.
\]
Then \(\mathbb ED_i=0\),
\[
\|D_i\|_2^2=\frac{1-q_i}{q_i},
\]
and the \(D_i\) are pairwise orthogonal. If \(c_i\) is the connection probability after forcing \(B_i\) open, then
\[
\langle\one_C,D_i\rangle=c_i-c,
\qquad
c_i\ge\tau_i.
\]
Bessel's inequality gives
\[
c(1-c)\ge\sum_iw_i(c_i-c)^2
\ge Q(\overline\tau-c)_+^2.
\]
Since \(c(1-c)\le1/4\),
\[
(\overline\tau-c)_+\le\frac1{2\sqrt Q},
\]
proving \eqref{eq:blockbessel}. \hfill\(\square\)

This shows that the rarity of individual all-open seeds is not itself fatal: enough disjoint seeds can make \(Q\) arbitrarily large.

\paragraph{9. Decision-tree deletion inequality.}
Gladkov's decision-tree framework permits certain configuration-dependent mixing operations while preserving the product law; see \href{https://arxiv.org/html/2408.08457v2}{[S4]}.

Let a decision tree explore a source configuration \(\omega\). On
\[
H=\{s\leftrightarrow A\},
\]
suppose it outputs \(a(\omega)\in A\) and an open \(s\)-to-\(a(\omega)\) path whose edges were queried. Define the closed trace
\[
Z(\omega)=\{e:\ e\text{ was queried and }\omega_e=0\}.
\]

\textbf{Lemma 9.1 (decision-tree deletion inequality).}
\[
\boxed{
\mathbb P_p(s\leftrightarrow t)
\ge
\mathbb E_\omega\left[
\one_H
\mathbb P_p^{\,G\setminus Z(\omega)}
\bigl(a(\omega)\leftrightarrow t\bigr)
\right].
}
\tag{9.1}\label{eq:decisiondelete}
\]

\emph{Proof sketch.}
Mix \(\omega\) with an independent \(\eta\), using \(\omega\) on queried edges and \(\eta\) on the rest, in the admissible decision-tree sense. The resulting configuration has the original product law. The queried open source path survives. Any \(\eta\)-open \(a(\omega)\)-to-\(t\) path that avoids \(Z(\omega)\) survives in the mixed configuration as well. Hence the mixed configuration connects \(s\) to \(t\) whenever the event on the right occurs. \hfill\(\square\)

If \(\Gamma_a\) is a canonical open \(a\)-to-\(t\) path and
\[
\pi_a(e)=\mathbb P_p(e\in\Gamma_a),
\]
then
\[
\mathbb P_p^{G\setminus Z}(a\leftrightarrow t)
\ge
\mathbb P_p(a\leftrightarrow t)-\sum_{e\in Z}\pi_a(e).
\]
Consequently,
\[
\boxed{
\begin{aligned}
\mathbb P_p(s\leftrightarrow t)\ge{}&
\mathbb P_p(s\leftrightarrow A)
\min_{a\in A}\mathbb P_p(a\leftrightarrow t)\\
&-
\mathbb E\left[
\one_H\sum_{e\in Z(\omega)}\pi_{a(\omega)}(e)
\right].
\end{aligned}
}
\tag{9.2}\label{eq:pathprofile}
\]
Thus queried open edges are harmless; the loss comes from target paths encountering source edges that were queried and found closed.

\paragraph{10. Closed-trace shell criterion suggested by the decision-tree reformulation.}
After a source-side history \(\mathcal H\), suppose one has an open explored source set, closed trace \(Z(\mathcal H)\), disjoint still-unqueried connector blocks \(B_i\), and associated shell contacts \(a_i\). Let
\[
q_i=\mathbb P(B_i\text{ all open}),
\quad
w_i=\frac{q_i}{1-q_i},
\quad
Q=\sum_iw_i,
\]
and define an averaged target-path profile
\[
\overline\pi(e)=\frac1Q\sum_iw_i\pi_{a_i}(e).
\]
A natural sufficient target is
\[
\boxed{
\mathfrak L
=
\mathbb E\left[
\one_{\mathrm{good}}\sum_{e\in Z}\overline\pi(e)
\right]\longrightarrow0.
}
\tag{CTS}\label{eq:cts}
\]
Equivalently, if \(I\) is sampled with probability \(w_i/Q\) and an independent canonical target path \(\Gamma_{a_I}\) is drawn, then
\[
\mathfrak L=\mathbb E|Z\cap\Gamma_{a_I}|.
\]
This led to the geometric formulation:
\[
\boxed{
\text{Find a source exploration and averaged target witnesses
whose expected closed-edge intersection tends to zero.}
}
\]
Two sufficient bounds were identified:
\[
\mathfrak L\le\rho L
\]
if closed-edge revealment is at most \(\rho\) and the averaged target path spends expected length at most \(L\) in the overlap region, and
\[
\sum_er(e)\overline\pi(e)
\le
\left(\sum_er(e)^2\right)^{1/2}
\left(\sum_e\overline\pi(e)^2\right)^{1/2}.
\]
This makes transparent why half-space or quasi-planar mechanisms can succeed: if target paths can be chosen in a region disjoint from the source's closed trace, the loss is zero.

\paragraph{11. Final audit: the proposed shell closure was invalid.}
This was the decisive correction.

Fix a shell configuration \(\xi\). The Kozma--Nitzan setup gives a deterministic set \(A_\xi\) of shell components satisfying
\[
\tau_a(\xi)
=
\mathbb P_{K_\xi}(a\leftrightarrow T)>1-\delta,
\qquad a\in A_\xi,
\]
and separately gives high probability that the source reaches \(A_\xi\).

After exposing a source-side history \(\mathcal H\), one may obtain fixed, still-unqueried connector blocks \(B_i\) leading to \(a_i\in A_\xi\). But the block-Bessel inequality conditioned on \(\xi,\mathcal H\) requires
\[
\tau_i(\xi,\mathcal H)
=
\mathbb P(a_i\leftrightarrow T\mid\xi,\mathcal H),
\]
not the pre-exploration probabilities \(\tau_{a_i}(\xi)\).

The valid conditional inequality has the form
\[
\boxed{
\mathbb P(o\leftrightarrow T\mid\xi,\mathcal H)
\ge
\frac1{Q(\mathcal H)}
\sum_iw_i\tau_i(\xi,\mathcal H)
-\frac1{2\sqrt{Q(\mathcal H)}}.
}
\tag{11.1}\label{eq:validconditional}
\]
The previously suggested replacement
\[
\mathbb P(o\leftrightarrow T\mid\xi,\mathcal H)
\stackrel{\mathrm{invalid}}{\ge}
\frac1{Q(\mathcal H)}
\sum_iw_i\tau_{a_i}(\xi)
-\frac1{2\sqrt{Q(\mathcal H)}}
\tag{11.2}\label{eq:invalidconditional}
\]
is unjustified.

The history \(\mathcal H\) contains negative information---edges queried and found closed. Those closed edges may form part of a cut between \(a_i\) and \(T\), so
\[
\tau_i(\xi,\mathcal H)
\]
can be much smaller than \(\tau_{a_i}(\xi)\), possibly zero.

The exact mismatch is:
\[
\begin{array}{c|c|c}
&\text{fixed connector blocks}&\text{high target probabilities}\\ \hline
\text{before exposing }\mathcal H&\text{not yet fixed}&\text{available}\\
\text{after exposing }\mathcal H&\text{available}&\text{not controlled}.
\end{array}
\]

The decision-tree form exposes the same gap:
\[
\mathbb P(o\leftrightarrow T)
\ge
\mathbb E\left[
\one_{\mathcal S}
\mathbb P^{K_\xi\setminus Z(\mathcal H)}
\bigl(a(\mathcal H)\leftrightarrow T\bigr)
\right].
\tag{11.3}
\]
To finish the proof one would need, for example,
\[
\boxed{
\mathbb E\left[
\one_{\mathcal S}
\left(
\mathbb P_{K_\xi}(a(\mathcal H)\leftrightarrow T)
-
\mathbb P^{K_\xi\setminus Z(\mathcal H)}
(a(\mathcal H)\leftrightarrow T)
\right)
\right]
=o(1).
}
\tag{11.4}\label{eq:missingclosedtrace}
\]
No proof of \eqref{eq:missingclosedtrace} was obtained. This is not a missing constant but the original gluing difficulty in another form.

\paragraph{12. Final list of concrete sufficient targets.}
The attempt therefore leaves several explicit ways a future proof could close:
\begin{enumerate}
\item Prove proportional local uniqueness:
\[
\theta(p)>0
\Longrightarrow
\lim_{K\to\infty}\limsup_{L\to\infty}\mathbb P_p(D_{L,K})=0.
\]

\item Prove the summable line lower-tail estimate \eqref{eq:linecriterion}.

\item Prove a finite-cluster radius-tail estimate strong enough for \eqref{eq:summableq}, for example
\[
q_{p_c}(R)=O(R^{-1-\alpha})
\]
for some \(\alpha>0\).

\item Prove the closed-trace robustness estimate \eqref{eq:missingclosedtrace}, or a sufficiently strong shell-specific variant.

\item Prove the Kozma--Nitzan high-probability gluing conjecture directly.
\end{enumerate}

\paragraph{Verification and outcome.}
The contact-mass lemma, line-gate lemma, radius-tail reduction, connector-block Bessel estimate, and decision-tree deletion formulation are retained as derived intermediate progress. They do not constitute a proof of the conjecture.

The final audit invalidates the proposed shell closure because target-goodness before exposing the source history does not imply target-goodness after conditioning on the closed edges revealed by that history.

\textbf{Outcome: substantial reductions and auxiliary inequalities; UNRESOLVED.}
No complete proof of
\[
\theta(p_c)=0
\]
is claimed.

\subsection{Sources}
\begin{itemize}
\item[S1] Rapha\"el Cerf, 2026 preprint on critical percolation and finite-cluster tails.
\url{https://arxiv.org/html/2608.23661v1}.
\textit{Used in the conversation for the 2026 status checkpoint and the finite-cluster-tail alternative.}

\item[S2] Gady Kozma and Asaf Nitzan, 2024 preprint on critical percolation and finite-graph connectivity inequalities.
\url{https://arxiv.org/html/2401.12397v1}.
\textit{Used for local uniqueness, shell constructions, the conjectural gluing inequality, and the line/half-space route.}

\item[S3] Thomas M. Liggett, Roberto H. Schonmann, and Alan M. Stacey, ``Domination by product measures,'' \emph{Annals of Probability} 25 (1997).
\url{https://projecteuclid.org/journals/annals-of-probability/volume-25/issue-1/Domination-by-product-measures/10.1214/aop/1024404279.pdf}.
\textit{Used for finite-range dependent coarse-graining.}

\item[S4] Alexey Gladkov, 2024 preprint on decision-tree inequalities for percolation connectivity.
\url{https://arxiv.org/html/2408.08457v2}.
\textit{Used for the decision-tree mixing viewpoint and the adaptive-switching obstruction.}

\item[S5] Classical supercritical finite-cluster tail reference cited during the conversation.
\url{https://www.jstor.org/stable/2244302}.
\textit{Background for the distinction between strictly supercritical tail estimates and estimates needed from the weaker hypothesis \(\theta(p)>0\).}

\item[S6] Slab/percolation reference cited during the final audit.
\url{https://arxiv.org/pdf/1401.7130}.
\textit{Used to illustrate how negative information from an exposed path can obstruct full-space gluing and how quasi-planar geometry can help in slab settings.}
\end{itemize}

\end{document}
